% chapters/ch22_epsilon_supply.tex
% Chapter 22: ε Supply Mechanism — ε as an Institutional Layer
% Source: NRMO_v5_updated.pdf, Chapter 22 (lines 9434-9500 of v5_layout.txt)

\chapter{$\varepsilon$ Supply Mechanism ---
$\varepsilon$ as an Institutional Layer}
\label{ch:epsilon-supply}

\nrmobox{Document status}{
\begin{tabular}{ll}
\textbf{Status}: & Foundational \\
\textbf{Layer}:  & READ ONLY (Immutable)
\end{tabular}\\[0.3em]
$\varepsilon$ is \textbf{not} a tuning parameter. It is a policy
artefact supplied by external governance. NRMO never generates,
modifies, or recommends $\varepsilon$.
}

% =====================================================================
\section{Basic Definition}
\label{sec:epsilon-supply-definition}

\begin{definition}[$\varepsilon$-Tolerance]
\label{def:epsilon-tolerance}
$\varepsilon \in (0, 1)$. $\varepsilon$ is the permissible risk
threshold supplied externally to NRMO.
\end{definition}

% =====================================================================
\section{Supply Entity}
\label{sec:epsilon-supply-entity}

$\varepsilon$ can be supplied only by the following principals:

\begin{enumerate}
\item \textbf{Human Governor} (human administrator) --- Final
$\varepsilon$ decision authority.
\item \textbf{Regulatory Constraint} --- Legal / institutional
upper limit.
\item \textbf{Domain Policy} --- Domain-specific constraints.
\item \textbf{Physical Bound} --- Upper limit by natural law.
\end{enumerate}

% =====================================================================
\section{Inviolable Properties}
\label{sec:epsilon-inviolable}

\nrmobox{$\varepsilon$ Inviolability Clause}{
\begin{itemize}
\item NRMO does \textbf{not} generate $\varepsilon$.
\item NRMO does \textbf{not} modify $\varepsilon$.
\item NRMO does \textbf{not} recommend $\varepsilon$.
\end{itemize}
\centering
These three prohibitions are the \textbf{hardest constraints in the
system}. Any mechanism that allows NRMO to self-modify $\varepsilon$
is a structural violation.
}

% =====================================================================
\section{Time Dependence}
\label{sec:epsilon-time-dependence}

$\varepsilon$ may depend on time and domain:

\begin{equation}
\varepsilon = \varepsilon(t, \mathrm{domain}).
\end{equation}

This dependence is set explicitly by the Human Governor. Automatic
time-variation of $\varepsilon$ by NRMO is prohibited.

% =====================================================================
\section{$\varepsilon$ Change Protocol (v2.1)}
\label{sec:epsilon-change-protocol}

The v2.1 MetaGovernanceLayer requires $\varepsilon$ changes to
follow this protocol:

\begin{enumerate}
\item Change request is issued only by Human Governor.
\item Change is recorded with signature in append-only log.
\item Old $\varepsilon$ value is not deleted (preserved permanently
for audit).
\item Reason for change is mandatory.
\end{enumerate}

\begin{lstlisting}[basicstyle=\ttfamily\small,frame=single]
epsilon_config = {
  'value':         float,        # Current epsilon value
  'ratified_date': timestamp,    # Last approval date
  'renewal_due':   timestamp,    # Renewal required by
  'ratified_by':   str,          # Governor identifier
  'domain':        str,          # Applicable domain
  'reason':        str,          # Reason for current value
  'audit_log': [                 # Append-only history
    {
      'old_value':  float,
      'new_value':  float,
      'changed_by': str,
      'reason':     str,
      'timestamp':  timestamp
    }
  ]
}
\end{lstlisting}

% =====================================================================
\section{Why $\varepsilon$ Must Be External}
\label{sec:epsilon-external-rationale}

If NRMO could self-modify $\varepsilon$:

\begin{itemize}
\item The system could relax its own constraints when under
pressure.
\item Self-referential $\varepsilon$ modification is a form of
corruption.
\item Human oversight would be structurally bypassed.
\end{itemize}

\noindent
This is why $\varepsilon$ externality is a constitutional
requirement, not merely a design choice.

% =====================================================================
\section{Open Questions (Future Work)}
\label{sec:epsilon-open-questions}

\begin{itemize}
\item Who has authority to approve $\varepsilon$ changes in
multi-stakeholder environments?
\item How are conflicting stakeholder preferences resolved?
\item What mechanisms ensure $\varepsilon$ reflects legitimate risk
tolerance?
\item Formalise democratic $\varepsilon$-governance protocol.
\end{itemize}

% =====================================================================
\section*{Connections to Other Parts}

\begin{itemize}
\item This chapter formalises Invariant~\ref{inv:gov-exec-sep}
($\varepsilon$ externally supplied) at the institutional layer.
\item Operational mechanism: Section~\ref{sec:parameter-renewal}
(parameter renewal protocol).
\item In Part~VII, the chance constraint $P_\pi(\tau_F \le T) \le
\varepsilon$ uses this institutional $\varepsilon$ as
exogenous input.
\item In Part~IX, $\Omega$ Full's adaptive tuning layer adjusts
\emph{behavioural} parameters within the constitutional
$\varepsilon$ ceiling but never the ceiling itself.
\end{itemize}
